\documentclass[10pt]{article}
\usepackage[a4paper,margin=20mm]{geometry}
\usepackage[T1]{fontenc}
\usepackage{lmodern,amsmath,amssymb,booktabs}
\usepackage[hidelinks]{hyperref}
\setcounter{MaxMatrixCols}{20}
\setlength{\arraycolsep}{3.5pt}
\setlength{\parindent}{0pt}
\setlength{\parskip}{4pt}
\allowdisplaybreaks
\numberwithin{equation}{section}
\title{Full Unicycle Model and Turning Control}
\author{}
\date{}
\begin{document}
\maketitle
\vspace{-2em}

\section{Conventions}
SI units; angles in radians; ideal rolling; $\cos\theta\ne0$.
\begin{equation}
\begin{aligned}
q&=(\psi,\theta,\phi,r,\gamma,x_G,y_G)^T,\\
\sigma&=(\sigma_1,\sigma_2,\sigma_3,\sigma_r,\sigma_g)^T,
\qquad x=(\sigma^T,q^T)^T,\qquad u=(F,M_2)^T.
\end{aligned}
\end{equation}
\begin{center}
\begin{tabular}{ll}
\toprule
Symbol & Meaning\\\midrule
$\psi,\theta,\phi$ & Heading, lean, wheel spin\\
$r,\gamma$ & Rod COM displacement along axle, body tilt about axle\\
$x_G,y_G$ & Wheel-center position in the global frame\\
$F=-F_1$ & Force on rod; opposite to notebook force on wheel\\
$M_2$ & Wheel drive torque; opposite reaction on body\\
\bottomrule
\end{tabular}
\end{center}
\begin{equation}
\begin{aligned}
\sigma_1&=\dot\theta,&\sigma_2&=\dot\phi+\dot\psi\sin\theta,&
\sigma_3&=\dot\psi\cos\theta,\\
\sigma_r&=\dot r-R\dot\theta,&\sigma_g&=\dot\gamma+\dot\psi\sin\theta.
\end{aligned}
\end{equation}
\textbf{Parameters.} $R$: wheel radius; $h$: wheel-center to body-COM distance;
$m_w,m_r,m_p$: wheel, rod, body masses. Inertias are about each component's COM:
$J_{W1}$ about the wheel axle, $J_{W2}$ about wheel $x,z$ axes,
$J_R$ about rod $x,z$ axes, and $J_{Px},J_{Py},J_{Pz}$ about body axes.
$B_R,B_P$ are rod and wheel--body viscous damping coefficients.
\begin{equation}
\begin{aligned}
(g,R,h)&=(9.81,0.253,0.025),& (m_w,m_r,m_p)&=(2.436,2.3,2.799),\\
(J_{W1},J_{W2},J_R)&=(0.09099921839,0.04591427768,0.0517629),\\
(J_{Px},J_{Py},J_{Pz})&=(0.01290418213,0.02090219895,0.01118711607),\\
(B_R,B_P)&=(0,0)\quad\text{(baseline)}.
\end{aligned}
\end{equation}

\section{Acceleration energy and pseudo forces}
\label{sec:appell-energy-forces}

\begin{equation}
\begin{gathered}
a=J_{Px}-J_{Pz}+m_ph^2,\qquad b=Rm_ph,\qquad c=J_{Pz}+J_R+J_{W2},\\
s_\gamma=\sin\gamma,\qquad c_\gamma=\cos\gamma,\qquad t_\theta=\tan\theta.
\end{gathered}
\end{equation}

\subsection*{Acceleration energy}
Up to terms independent of $\dot\sigma$,
\begin{align}
\mathcal S={}&\tfrac12\left[c+R^2(m_p+m_w)+2bc_\gamma+m_rr^2+ac_\gamma^2\right]\dot\sigma_1^2\\
&+\tfrac12\left[J_{W1}+R^2(m_p+m_r+m_w)\right]\dot\sigma_2^2
+\tfrac12\left[c+m_rr^2+as_\gamma^2\right]\dot\sigma_3^2\\
&+\tfrac12m_r\dot\sigma_r^2
+\tfrac12(J_{Py}+m_ph^2)\dot\sigma_g^2\\
&-(b+ac_\gamma)s_\gamma\dot\sigma_1\dot\sigma_3
-Rm_rr\dot\sigma_2\dot\sigma_3
+bc_\gamma\dot\sigma_2\dot\sigma_g\notag\\
&+\dot\sigma_1\Bigl\{\begin{aligned}[t]
&Rm_rr\sigma_1^2+2m_rr\sigma_1\sigma_r\\
&+(b+ac_\gamma)s_\gamma t_\theta\sigma_1\sigma_3
-2(b+ac_\gamma)s_\gamma\sigma_1\sigma_g\\
&-\left[J_{W1}+R^2(m_p+m_w)+bc_\gamma+Rm_rrt_\theta\right]\sigma_2\sigma_3\\
&+\left[bc_\gamma+m_rr^2+ac_\gamma^2+c\right]t_\theta\sigma_3^2\\
&+\left[J_{Px}-J_{Py}-J_{Pz}-2bc_\gamma-2ac_\gamma^2\right]\sigma_3\sigma_g
\end{aligned}\Bigr\}\notag\\
&+\dot\sigma_2\Bigl\{\begin{aligned}[t]
&-bs_\gamma(\sigma_3^2+\sigma_g^2)-2Rm_r\sigma_3\sigma_r\\
&+\left[R^2(m_p-m_r+m_w)+bc_\gamma+Rm_rrt_\theta\right]\sigma_1\sigma_3
\end{aligned}\Bigr\}\notag\\
&+\dot\sigma_3\Bigl\{\begin{aligned}[t]
&J_{W1}\sigma_1\sigma_2+bs_\gamma\sigma_2\sigma_3
+2m_rr\sigma_3\sigma_r\\
&+\left[Rm_rr-(m_rr^2+as_\gamma^2+c)t_\theta\right]\sigma_1\sigma_3\\
&+\left[-J_{Px}+J_{Py}+J_{Pz}+2as_\gamma^2\right]\sigma_1\sigma_g\\
&-as_\gamma c_\gamma t_\theta\sigma_3^2
+2as_\gamma c_\gamma\sigma_3\sigma_g
\end{aligned}\Bigr\}\notag\\
&+\dot\sigma_r\left[Rm_r\sigma_2\sigma_3-m_rr(\sigma_1^2+\sigma_3^2)\right]\notag\\
&+\dot\sigma_g\Bigl\{\begin{aligned}[t]
&bs_\gamma t_\theta\sigma_2\sigma_3
+(b+ac_\gamma)s_\gamma\sigma_1^2\\
&+(a+bc_\gamma-2as_\gamma^2)\sigma_1\sigma_3
-as_\gamma c_\gamma\sigma_3^2
\end{aligned}\Bigr\}.
\end{align}

\subsection*{Pseudo forces}
In the order $(\sigma_1,\sigma_2,\sigma_3,\sigma_r,\sigma_g)$,
\begin{equation}
\boxed{\begin{aligned}
\mathcal Q_1={}&RF-B_RR(R\sigma_1+\sigma_r)
+g[R(m_p+m_w)+hm_p\cos\gamma]\sin\theta-gm_rr\cos\theta,\\
\mathcal Q_2={}&M_2-B_P(\sigma_2-\sigma_g),\\
\mathcal Q_3={}&-ghm_p\sin\gamma\sin\theta,\\
\mathcal Q_r={}&F-B_R(R\sigma_1+\sigma_r)-gm_r\sin\theta,\\
\mathcal Q_g={}&-M_2+B_P(\sigma_2-\sigma_g)+ghm_p\sin\gamma\cos\theta.
\end{aligned}}
\end{equation}

At fixed $(q,\sigma)$, apply
\begin{equation}
\frac{\partial\mathcal S}{\partial\dot\sigma_i}=\mathcal Q_i,
\qquad i\in\{1,2,3,r,g\}.
\end{equation}

Collecting all five equations in the same way, the coefficients of
$\dot\sigma$ form $M$, the remaining gravity and inertial terms form $n$,
and the damping and actuator terms form $D\sigma$ and $Q_u u$:
\begin{equation}
\boxed{M(q)\dot\sigma+n(q,\sigma)+D\sigma=Q_u u,
\qquad \dot q=k(q,\sigma).}
\end{equation}


\section{Full nonlinear model}
\begin{equation}
\boxed{M(q)\dot\sigma+n(q,\sigma)+D\sigma=Q_u u,\qquad \dot q=k(q,\sigma).}
\end{equation}
\begin{equation}
Q_u=\begin{bmatrix}R&0\\0&1\\0&0\\1&0\\0&-1\end{bmatrix},\qquad
D=\begin{bmatrix}
B_RR^2&0&0&B_RR&0\\0&B_P&0&0&-B_P\\0&0&0&0&0\\
B_RR&0&0&B_R&0\\0&-B_P&0&0&B_P
\end{bmatrix}.
\end{equation}
The residual notation used in the turning calculations is
\begin{equation}
\mathcal R(x,u)=n(q,\sigma)+D\sigma-Q_u u,\qquad
f(x,u)=\begin{bmatrix}-M(q)^{-1}\mathcal R(x,u)\\k(q,\sigma)\end{bmatrix}.
\end{equation}
Here $\mathcal R$ is the vector denoted by $C(x,u)$ in the earlier notes;
$D$ remains the damping matrix and $Q_u$ the input matrix.
\subsection*{Mass matrix}
\begin{equation}a=J_{Px}-J_{Pz}+m_ph^2,\qquad b=Rm_ph,\qquad c=J_{Pz}+J_R+J_{W2},\end{equation}

\begin{equation}s_\gamma=\sin\gamma,\quad c_\gamma=\cos\gamma,\quad t_\theta=\tan\theta.\end{equation}

\begin{equation}
\boxed{M=\begin{bmatrix}
M_{11}&0&M_{13}&0&0\\
0&M_{22}&-Rm_rr&0&bc_\gamma\\
M_{13}&-Rm_rr&M_{33}&0&0\\
0&0&0&m_r&0\\
0&bc_\gamma&0&0&J_{Py}+m_ph^2
\end{bmatrix}}
\end{equation}

\begin{equation}
\begin{aligned}
M_{11}&=c+R^2(m_p+m_w)+2bc_\gamma+m_rr^2+ac_\gamma^2,\\
M_{13}&=-(b+ac_\gamma)s_\gamma,\\
M_{22}&=J_{W1}+R^2(m_p+m_r+m_w),\\
M_{33}&=c+m_rr^2+as_\gamma^2.
\end{aligned}
\end{equation}
\subsection*{Gravity and inertial terms}
\begin{equation}
\begin{aligned}
n_1={}&-g\left[R(m_p+m_w)+hm_pc_\gamma\right]\sin\theta
       +gm_rr\cos\theta\\
&+Rm_rr\sigma_1^2+2m_rr\sigma_1\sigma_r\\
&+(b+ac_\gamma)s_\gamma t_\theta\sigma_1\sigma_3
 -2(b+ac_\gamma)s_\gamma\sigma_1\sigma_g\\
&-\left[J_{W1}+R^2(m_p+m_w)+bc_\gamma+Rm_rrt_\theta\right]\sigma_2\sigma_3\\
&+\left[bc_\gamma+m_rr^2+ac_\gamma^2+c\right]t_\theta\sigma_3^2\\
&+\left[J_{Px}-J_{Py}-J_{Pz}-2bc_\gamma-2ac_\gamma^2\right]\sigma_3\sigma_g.
\end{aligned}
\end{equation}

\begin{equation}
\begin{aligned}
n_2={}&-bs_\gamma(\sigma_3^2+\sigma_g^2)-2Rm_r\sigma_3\sigma_r\\
&+\left[R^2(m_p-m_r+m_w)+bc_\gamma+Rm_rrt_\theta\right]\sigma_1\sigma_3.
\end{aligned}
\end{equation}

\begin{equation}
\begin{aligned}
n_3={}&J_{W1}\sigma_1\sigma_2+bs_\gamma\sigma_2\sigma_3
       +ghm_ps_\gamma\sin\theta+2m_rr\sigma_3\sigma_r\\
&+\left[Rm_rr-(m_rr^2+as_\gamma^2+c)t_\theta\right]\sigma_1\sigma_3\\
&+\left[-J_{Px}+J_{Py}+J_{Pz}+2as_\gamma^2\right]\sigma_1\sigma_g\\
&-as_\gamma c_\gamma t_\theta\sigma_3^2
 +2as_\gamma c_\gamma\sigma_3\sigma_g.
\end{aligned}
\end{equation}

\begin{equation}n_4=Rm_r\sigma_2\sigma_3+gm_r\sin\theta-m_rr(\sigma_1^2+\sigma_3^2).\end{equation}

\begin{equation}
\begin{aligned}
n_5={}&bs_\gamma t_\theta\sigma_2\sigma_3-ghm_ps_\gamma\cos\theta\\
&+(b+ac_\gamma)s_\gamma\sigma_1^2
 +(a+bc_\gamma-2as_\gamma^2)\sigma_1\sigma_3
 -as_\gamma c_\gamma\sigma_3^2.
\end{aligned}
\end{equation}
\subsection*{Kinematics}
\begin{equation}
\boxed{
\begin{aligned}
\dot\psi&=\frac{\sigma_3}{\cos\theta},\\
\dot\theta&=\sigma_1,\\
\dot\phi&=\sigma_2-\sigma_3\tan\theta,\\
\dot r&=R\sigma_1+\sigma_r,\\
\dot\gamma&=\sigma_g-\sigma_3\tan\theta,\\
\dot x_G&=R\sigma_1\sin\psi\cos\theta+R\sigma_2\cos\psi,\\
\dot y_G&=-R\sigma_1\cos\psi\cos\theta+R\sigma_2\sin\psi.
\end{aligned}}
\end{equation}

\section{Linearized model}
Upright straight rolling along $+x$: $V=v/R$ is constant,
$\sigma_e=(0,V,0,0,0)^T$, $\theta_e=r_e=\gamma_e=\psi_e=0$,
$\dot\phi_e=V$, $\dot x_{Ge}=RV$, $u_e=(0,B_PV)^T$.
All variables below are perturbations; $V=0$ gives the rest linearization.
\begin{equation}
\begin{aligned}
I_\ell&=J_{Px}+J_R+J_{W2}+R^2(m_p+m_w)+2b+m_ph^2,\\
G_\ell&=g[R(m_p+m_w)+m_ph],& I_w^*&=J_{W1}+R^2(m_p+m_r+m_w),\\
I_g&=J_{Py}+m_ph^2,&G_g&=gm_ph,\\
L&=J_{W1}+R^2(m_p+m_w)+b,&\kappa&=J_{W1}V/c.
\end{aligned}
\end{equation}
\begin{equation}
\begin{aligned}
I_\ell\dot\sigma_1&=G_\ell\theta-gm_r r+LV\sigma_3-B_RR(R\sigma_1+\sigma_r)+RF,\\
I_w^*\dot\sigma_2+b\dot\sigma_g&=-B_P(\sigma_2-\sigma_g)+M_2,\\
c\dot\sigma_3&=-J_{W1}V\sigma_1,\\
\dot\sigma_r&=-g\theta-RV\sigma_3-\frac{B_R}{m_r}(R\sigma_1+\sigma_r)+\frac{F}{m_r},\\
b\dot\sigma_2+I_g\dot\sigma_g&=G_g\gamma+B_P(\sigma_2-\sigma_g)-M_2.
\end{aligned}
\end{equation}
\begin{equation}
\begin{aligned}
\dot\psi&=\sigma_3,&\dot\theta&=\sigma_1,&\dot\phi&=\sigma_2,\\
\dot r&=R\sigma_1+\sigma_r,&\dot\gamma&=\sigma_g,\\
\dot x_G&=R\sigma_2,&\dot y_G&=-R\sigma_1+RV\psi.
\end{aligned}
\end{equation}
These equations specify $\delta\dot x=A(V)\delta x+B\delta u$ in the stated 12-state order.

\section{Lean--yaw-rate invariant}
At the upright straight-rolling reference, $n_{3e}=0$. Taking partial
 derivatives with respect to the states gives
\begin{equation}
\delta n_3=\left.\frac{\partial n_3}{\partial x}\right|_e\delta x
=J_{W1}V\,\delta\sigma_1,
\end{equation}
since only $\partial n_3/\partial\sigma_1$ is nonzero at this reference.
This is a first-order variation, not a time derivative. The third dynamic
row and lean kinematics give
\begin{equation}
c\,\delta\dot\sigma_3+J_{W1}V\,\delta\sigma_1=0,\qquad
\delta\dot\theta=\delta\sigma_1,\qquad
\frac{d}{dt}(\delta\sigma_3+\kappa\delta\theta)=0.
\end{equation}
Suppressing perturbation symbols, for the fixed-speed linear model only:
\begin{equation}
\boxed{C=\sigma_3+\kappa\theta=\text{constant},\qquad
\dot\psi=C-\kappa\theta.}
\end{equation}
\begin{equation}
w^T=(0,0,1,0,0,0,\kappa,0,0,0,0,0),\qquad
w^TA=w^TB=w^T(A-BK)=0.
\end{equation}
\begin{equation}
\sigma_3(0)=0\ \Longrightarrow\ C=\kappa\theta_0,\qquad
\sigma_{3\infty}+\kappa\theta_\infty=C.
\end{equation}
This relates lean to yaw \emph{rate}, not heading. For $V\ne0$, zero terminal yaw rate requires
$\theta_\infty=\theta_0$ for the initial condition above. It is not an exact nonlinear conservation law.

\section{Equilibrium manifolds}
\subsection*{General relative equilibrium}
Let $V=\dot\phi_e$ and $\Omega=\dot\psi_e$ be constant, with fixed $\theta_e,r_e,\gamma_e$.
\begin{equation}
\sigma_e=(0,V+\Omega\sin\theta_e,\Omega\cos\theta_e,0,\Omega\sin\theta_e)^T.
\end{equation}
Writing $n_{ie}=n_i(q_e,\sigma_e)$, the complete dynamic conditions are
\begin{equation}
\boxed{n_{1e}-Rn_{4e}=0,\quad n_{3e}=0,\quad n_{2e}+n_{5e}=0,\qquad
F_e=n_{4e},\quad M_{2e}=n_{2e}+B_PV.}
\end{equation}
A stationary equilibrium additionally has $V=\Omega=0$; rolling and turning references generally have $\dot x_e\ne0$.

\subsection*{Upright-body branch: $\gamma_e=0$}
Suppressing equilibrium subscripts, define
\begin{equation}
G_T=g[R(m_p+m_w+m_r)+hm_p],\quad
H=J_{W1}+R^2(m_p+m_w+m_r)+b,\quad
K_T=c-J_{Py}-m_ph^2-b-H.
\end{equation}
\begin{equation}
\begin{aligned}
0={}&-G_T\sin\theta+gm_r r\cos\theta\\
&-V\Omega(H\cos\theta+Rm_r r\sin\theta)\\
&+\Omega^2[(K_T+m_r r^2)\sin\theta\cos\theta+Rm_r r\cos(2\theta)].
\end{aligned}
\end{equation}
\begin{equation}
\begin{aligned}
F_e&=m_r[R(V+\Omega\sin\theta)\Omega\cos\theta+g\sin\theta-r\Omega^2\cos^2\theta],\\
M_{2e}&=B_PV.
\end{aligned}
\end{equation}
\textbf{Straight rolling} ($\Omega=0$; stationary if also $V=0$):
\begin{equation}
\boxed{r_e=\frac{G_T}{gm_r}\tan\theta_e,\qquad
F_e=m_rg\sin\theta_e,\qquad M_{2e}=B_PV.}
\end{equation}
\textbf{Upright turning} ($\theta_e=0$):
\begin{equation}
\boxed{r_e=\frac{HV\Omega_e}{m_r(g+R\Omega_e^2)},\qquad
F_e=m_r(RV\Omega_e-r_e\Omega_e^2),\qquad M_{2e}=B_PV.}
\end{equation}
\textbf{Local manifold tangent} at $(\theta,r,\Omega)=(0,0,0)$ for fixed $V$:
\begin{equation}
gm_r r-G_T\theta-HV\Omega=0.
\end{equation}


\section{General turning manifold control}
\subsection{Steady-turning residual equations}
Use the $\gamma_e=0$ branch. The wheel-phase rate $\nu$ in the earlier
notes is denoted by $V=\dot\phi_e=v/R$ here; $\Omega=\dot\psi_e$ is the yaw rate.
The equilibrium shape and pseudovelocities are
\begin{equation}
\begin{aligned}
\sigma_e&=(0,V+\Omega\sin\theta_e,\Omega\cos\theta_e,0,\Omega\sin\theta_e)^T,\\
q_e(t)&=(\psi_e+\Omega t,\theta_e,\phi_e+Vt,r_e,0,x_{Ge}(t),y_{Ge}(t))^T.
\end{aligned}
\end{equation}
The wheel-center coordinates follow the kinematics. With
$\mathcal R_{ie}=\mathcal R_i(x_e,u_e)$ and the previously defined
$a=J_{Px}-J_{Pz}+m_ph^2$, $b=Rhm_p$, $c=J_{Pz}+J_R+J_{W2}$,
the five residual equations in the supplied calculations become
\begin{align}
0=\mathcal R_{1e}={}&-RF_e-g[R(m_p+m_w)+hm_p]\sin\theta_e
+gm_r r_e\cos\theta_e\notag\\
&-\sigma_{2e}\sigma_{3e}
[J_{W1}+R^2(m_p+m_w)+b+Rm_r r_e\tan\theta_e]\notag\\
&+\sigma_{3e}^2(b+m_r r_e^2+a+c)\tan\theta_e\notag\\
&+\sigma_{3e}\sigma_{ge}(J_{Px}-J_{Py}-J_{Pz}-2b-2a),\\
0=\mathcal R_{2e}={}&B_P(\sigma_{2e}-\sigma_{ge})-M_{2e},\\
0=\mathcal R_{3e}={}&0,\\
0=\mathcal R_{4e}={}&-F_e+Rm_r\sigma_{2e}\sigma_{3e}
+gm_r\sin\theta_e-m_r r_e\sigma_{3e}^2,\\
0=\mathcal R_{5e}={}&-\mathcal R_{2e}.
\end{align}
Thus $\mathcal R_{4e}=0$ gives $F_e$ as a function of $r_e$;
$\mathcal R_{2e}=0$ gives $M_{2e}=B_PV$. Substituting these into
$\mathcal R_{1e}=0$ determines $r_e$. Select the desired continuous solution
branch and verify all five residuals:
\begin{equation}
\boxed{\mathcal R(x_e,u_e)=0,\qquad
f(x_e,u_e)=\begin{bmatrix}0_{5\times1}\\k(q_e,\sigma_e)\end{bmatrix}.}
\end{equation}
The identity $\mathcal R_{3e}=0$ holds on this equilibrium branch;
its partial derivatives need not vanish for perturbations away from the branch.

\subsection{Nine-state linearization}
Remove wheel phase and global position from the feedback state:
\begin{equation}
y=x_I=(\sigma_1,\sigma_2,\sigma_3,\sigma_r,\sigma_g,\psi,\theta,r,\gamma)^T,
\qquad I=(1,2,3,4,5,6,7,9,10).
\end{equation}
Let $y_r(t)$ retain the equilibrium shape and rates, with
$\psi_r(t)=\psi_e+\Omega t$. For $e=y-y_r(t)$ and $\delta u=u-u_e$,
\begin{equation}
\dot e=A_e e+B_e\delta u+O((\|e\|+\|\delta u\|)^2),\qquad
A_e=\left.\frac{\partial f_I}{\partial y}\right|_e,\quad
B_e=\left.\frac{\partial f_I}{\partial u}\right|_e.
\end{equation}
Since $\mathcal R_e=0$, differentiation of $-M^{-1}\mathcal R$ yields
\begin{equation}
\boxed{A_e=\begin{bmatrix}-M_e^{-1}\mathcal R_y|_e\\H_e\end{bmatrix}
\in\mathbb R^{9\times9},\qquad
B_e=\begin{bmatrix}M_e^{-1}Q_u\\0_{4\times2}\end{bmatrix}
\in\mathbb R^{9\times2}.}
\end{equation}
The derivatives in $\mathcal R_y$ are evaluated at fixed input $u_e$;
they are not derivatives along the equilibrium manifold.
\begin{equation}
\mathcal R_y|_e=
[\mathcal R_{\sigma_1}\ \mathcal R_{\sigma_2}\ \mathcal R_{\sigma_3}\
\mathcal R_{\sigma_r}\ \mathcal R_{\sigma_g}\ 0_{5\times1}\
\mathcal R_\theta\ \mathcal R_r\ \mathcal R_\gamma]_e.
\end{equation}
The kinematic block is
\begin{equation}
H_e=\begin{bmatrix}
0&0&\sec\theta_e&0&0&0&\sigma_{3e}\sec\theta_e\tan\theta_e&0&0\\
1&0&0&0&0&0&0&0&0\\
R&0&0&1&0&0&0&0&0\\
0&0&-\tan\theta_e&0&1&0&-\sigma_{3e}\sec^2\theta_e&0&0
\end{bmatrix}.
\end{equation}
These matrices are constant along the reference because $f_I$ is independent
of heading and the omitted coordinates.

\subsection{Local target selection}
At each candidate equilibrium form the controllability matrix
\begin{equation}
\mathcal C_e=[B_e\ A_eB_e\ \cdots\ A_e^8B_e]\in\mathbb R^{9\times18}.
\end{equation}
The supplied design uses a branch with numerical rank eight. Let the unit
vector $w_e$ span its left nullspace:
\begin{equation}
w_e^T\mathcal C_e=0,\qquad \|w_e\|_2=1.
\end{equation}
For prescribed $\theta_e,V$, choose $\Omega$ by solving
\begin{equation}
\boxed{h(\Omega)=w_e(\Omega)^T[y_0-y_e(\theta_e,\Omega,V)]=0,}
\end{equation}
where $y_e=y_r(0)$. Recompute the equilibrium and Jacobians for each candidate
and keep the sign of $w_e$ consistent during root finding. This matches the
initial error to the local controllable subspace; it is not an exact nonlinear
conservation law or a guarantee of nonlinear reachability.

\subsection{SVD basis and transverse LQR}
Take the full singular value decomposition
\begin{equation}
\mathcal C_e=U\Sigma W^T,\qquad
\mathcal C_e\mathcal C_e^T u_i=s_i^2u_i,\qquad
s_1\ge\cdots\ge s_9\ge0.
\end{equation}
The left singular vectors $u_i$ form an orthonormal basis. For numerical rank eight,
\begin{equation}
T=[u_1,\ldots,u_8]\in\mathbb R^{9\times8},\qquad w_e=u_9,\qquad
T^TT=I_8,\quad T^Tw_e=0,\quad e=TT^Te+w_ew_e^Te.
\end{equation}
The reported implementation counts $s_i>10^{-9}s_1$. Rank depends on scaling
and tolerance; recompute the basis at every new target.
For $z=T^Te$, use
\begin{equation}
A_c=T^TA_eT,\qquad B_c=T^TB_e,\qquad Q_c=T^TQT,
\end{equation}
\begin{equation}
Q=\operatorname{diag}(2,1,1,2,1,0.5,20,20,5),\qquad
R_u=\operatorname{diag}(0.1,0.1).
\end{equation}
The reduced LQR gain and lifted feedback gain are
\begin{equation}
K_c=R_u^{-1}B_c^TP,\qquad K=K_cT^T,\qquad
A_c^TP+PA_c-PB_cR_u^{-1}B_c^TP+Q_c=0,
\end{equation}
where $P$ is the stabilizing Riccati solution. The projected model applies on
$w_e^Te=0$; otherwise its dynamics also include
$T^TA_ew_e(w_e^Te)$. The unactuated component must therefore be matched in
reference selection rather than discarded.

\subsection{Feedback implementation}
Wrap the heading error and apply componentwise actuator saturation:
\begin{equation}
e_6\leftarrow\operatorname{atan2}(\sin e_6,\cos e_6),\qquad
\boxed{u(t)=\operatorname{sat}_{[-u_{\max},u_{\max}]}(u_e-Ke(t)),}
\end{equation}
\begin{equation}
u_{\max}=\begin{bmatrix}30\,\mathrm N\\15\,\mathrm{N\,m}\end{bmatrix}.
\end{equation}
The target and gain remain fixed during each stage. Wheel phase and global
position are excluded from feedback. Local transverse decay rates are given
by the eigenvalues of $A_c-B_cK_c$, not by the singular values of $\mathcal C_e$.
These local unsaturated rates do not establish nonlinear stability under saturation.

\subsection*{Calculation sequence}
\begin{enumerate}
\item Specify $V$, $\theta_e$, and the initial state $y_0$.
\item For each candidate $\Omega$, solve for $r_e,F_e,M_{2e}$ and check
$\mathcal R_e=0$.
\item Form $A_e,B_e$, check rank eight, and compute $w_e$.
\item Solve $h(\Omega)=0$; construct $T$ and the LQR gain at the selected target.
\end{enumerate}


\par\medskip
\noindent\rule{\linewidth}{0.2pt}

\textbf{Below we assume fully linearized system and design controllers}

\noindent\rule{\linewidth}{0.2pt}
\par\medskip

\section{Fully linearized lateral reduced controller}

Consider the first-order model about upright straight rolling
at fixed wheel-phase rate $V=v/R\ne0$.
The lateral and longitudinal subsystems decouple to first order
at this reference; speed-dependent lean--yaw coupling remains within
the lateral subsystem. The lateral input is the rod force $F$.
All equations in this section refer to this fixed linear model;
the turning target is a relative equilibrium of the linear model.

\subsection{Lateral dynamics and invariant}

The lateral state vector is ordered as
\begin{equation}
x_{\rm lat}
=(\sigma_1,\sigma_r,r,\theta,\sigma_3,\psi)^T,
\end{equation}
with the linearized kinematic relations
\begin{equation}
\dot\theta=\sigma_1,\qquad
\dot r=R\sigma_1+\sigma_r,\qquad
\dot\psi=\sigma_3.
\end{equation}

The six-state lateral model is
\begin{equation}
\begin{aligned}
I_\ell\dot\sigma_1
  &=G_\ell\theta-gm_r r+LV\sigma_3
    -B_RR(R\sigma_1+\sigma_r)+RF,\\
\dot\sigma_r
  &=-g\theta-RV\sigma_3
    -\frac{B_R}{m_r}(R\sigma_1+\sigma_r)
    +\frac{F}{m_r},\\
\dot r&=R\sigma_1+\sigma_r,\\
\dot\theta&=\sigma_1,\\
\dot\sigma_3&=-\kappa\sigma_1,
\qquad \kappa=\frac{J_{W1}V}{c},\\
\dot\psi&=\sigma_3.
\end{aligned}
\end{equation}

\subsection{Lateral model and SVD reduction}

Use the lateral state ordering
\begin{equation}
x_{\rm lat}
=(\sigma_1,\sigma_r,r,\theta,\sigma_3,\psi)^T.
\end{equation}
At the upright straight-rolling reference with fixed $V$,
the linearized lateral dynamics are
\begin{equation}
\dot x_{\rm lat}=A_{\rm lat}x_{\rm lat}+B_{\rm lat}F,
\end{equation}
where
\begin{equation}
A_{\rm lat}=
\begin{bmatrix}
-\dfrac{B_RR^2}{I_\ell}
&-\dfrac{B_RR}{I_\ell}
&-\dfrac{gm_r}{I_\ell}
&\dfrac{G_\ell}{I_\ell}
&\dfrac{LV}{I_\ell}&0\\
-\dfrac{B_RR}{m_r}
&-\dfrac{B_R}{m_r}
&0&-g&-RV&0\\
R&1&0&0&0&0\\
1&0&0&0&0&0\\
-\kappa&0&0&0&0&0\\
0&0&0&0&1&0
\end{bmatrix},
\qquad
B_{\rm lat}=
\begin{bmatrix}
R/I_\ell\\
1/m_r\\
0\\
0\\
0\\
0
\end{bmatrix},
\qquad
\kappa=\frac{J_{W1}V}{c}.
\end{equation}

\subsection{Constructing the left-null vector and identifying it with SVD}

For this fixed linearization, set $A_e=A_{\rm lat}(V)$ and
$B_e=B_{\rm lat}$. Form
\begin{equation}
\mathcal C_e=[B_e\ A_eB_e\ \cdots\ A_e^5B_e]
=U_e\Sigma_eW_e^T.
\end{equation}
Assume that the selected operating speed gives
$\operatorname{rank}(\mathcal C_e)=5$. The SVD then supplies
\begin{equation}
T_{\rm svd}=[u_1,\ldots,u_5],\qquad
w_{\rm svd}=u_6,\qquad
\ker(\mathcal C_e^T)=\operatorname{span}\{w_{\rm svd}\}.
\end{equation}

We first construct an analytic candidate from the model equations,
rather than assuming an expression for the singular vector:
\begin{equation}
\widetilde w=(0,0,0,\kappa,1,0)^T,
\qquad
\widetilde w^Ty=\kappa\theta+\sigma_3,
\qquad y=x_{\rm lat}.
\end{equation}
Since $\kappa$ is constant and
$\dot\sigma_3+\kappa\dot\theta=0$, direct multiplication gives
\begin{equation}
\widetilde w^TA_e=0,\qquad \widetilde w^TB_e=0,
\qquad \widetilde w^T\mathcal C_e=0.
\end{equation}
Normalize this candidate by defining
\begin{equation}
\eta=\sqrt{1+\kappa^2},\qquad
\boxed{w_e=\frac{\widetilde w}{\eta}
=\frac{1}{\eta}(0,0,0,\kappa,1,0)^T.}
\end{equation}

Since $\widetilde w^T\mathcal C_e=0$ should only have one solution $\widetilde w_e$ direction for the assumed rank-four controllability, this analytic candidate must be correct up to a sign. 

\subsection{Explicit initial state and target selection}

Fix $V$, the desired lean $\theta_e$, and the initial reference heading
$\psi_{e0}$ (for example, $\psi_{e0}=\psi(0)$).
For each candidate yaw rate $\Omega$, define
\begin{equation}
G_T=G_\ell+Rm_rg,\qquad H=L+R^2m_r,\qquad
r_e(\Omega)=\frac{G_T\theta_e+HV\Omega}{gm_r},\qquad
F_e(\Omega)=m_r(g\theta_e+RV\Omega).
\end{equation}
The initial state and the initial point of the reference trajectory are
\begin{equation}
y_0=\begin{bmatrix}
\sigma_1(0)\\\sigma_r(0)\\r(0)\\\theta(0)\\\sigma_3(0)\\\psi(0)
\end{bmatrix},\qquad
y_e(\theta_e,\Omega,V)=\begin{bmatrix}
0\\0\\r_e(\Omega)\\\theta_e\\\Omega\\\psi_{e0}
\end{bmatrix}.
\end{equation}
Thus, explicitly,
\begin{equation}
\boxed{y_0-y_e=\begin{bmatrix}
\sigma_1(0)\\
\sigma_r(0)\\
r(0)-\dfrac{G_T\theta_e+HV\Omega}{gm_r}\\
\theta(0)-\theta_e\\
\sigma_3(0)-\Omega\\
\psi(0)-\psi_{e0}
\end{bmatrix}.}
\end{equation}
Choose $\Omega$ by solving the SVD compatibility condition
\begin{equation}
\boxed{h(\Omega)=w_e(\Omega)^T
[y_0-y_e(\theta_e,\Omega,V)]=0.}
\end{equation}
Here the notation $w_e(\Omega)$ follows the general turning design.
For the present fixed straight-rolling linearization, however,
$A_e$, $B_e$, and $w_e$ are independent of the candidate $\Omega$.
Substitution therefore gives
\begin{equation}
h(\Omega)=\frac{\kappa[\theta(0)-\theta_e]+\sigma_3(0)-\Omega}{\eta},
\qquad
\boxed{\Omega_e=\sigma_3(0)+\kappa[\theta(0)-\theta_e].}
\end{equation}
The other initial components do not enter this compatibility condition
because their coefficients in $w_e$ vanish; they still enter the feedback.

At the selected root, let $r_e=r_e(\Omega_e)$ and $F_e=F_e(\Omega_e)$.
Define
\begin{equation}
\psi_r(t)=\psi_{e0}+\Omega_e t,\qquad
y_r(t)=\begin{bmatrix}0&0&r_e&\theta_e&\Omega_e&\psi_r(t)\end{bmatrix}^T,
\qquad e=y-y_r(t),\qquad \delta F=F-F_e.
\end{equation}
This is a relative equilibrium, not a constant state when $\Omega_e\ne0$.
It satisfies the linear dynamics, so
\begin{equation}
\dot e=A_e e+B_e\delta F,\qquad
w_e^Te(t)=w_e^Te(0)=0.
\end{equation}

\subsection{An analytic orthonormal basis for the controllable subspace}

An explicit orthonormal complement of $w_e$ is
\begin{equation}
\boxed{
T=\begin{bmatrix}
1&0&0&0&0\\
0&1&0&0&0\\
0&0&1&0&0\\
0&0&0&1/\eta&0\\
0&0&0&-\kappa/\eta&0\\
0&0&0&0&1
\end{bmatrix}.}
\end{equation}
Direct calculation verifies
\begin{equation}
T^TT=I_5,\qquad T^Tw_e=0,\qquad
TT^T=I_6-w_ew_e^T=T_{\rm svd}T_{\rm svd}^T.
\end{equation}
Thus $T$ spans the same controllable subspace as $T_{\rm svd}$.
It is an analytic choice of basis, not generally the individual left
singular vectors. There exists an orthogonal $O\in\mathbb R^{5\times5}$
such that $T=T_{\rm svd}O$. Using either basis with consistently
transformed cost weights gives the same lifted LQR feedback.

The projected errors are explicitly
\begin{equation}
\boxed{z=T^Te=\begin{bmatrix}
\sigma_1\\\sigma_r\\r-r_e\\
\bigl[(\theta-\theta_e)-\kappa(\sigma_3-\Omega_e)\bigr]/\eta\\
\psi-\psi_r(t)
\end{bmatrix}.}
\end{equation}
On the compatible subspace,
\begin{equation}
\sigma_3-\Omega_e=-\kappa(\theta-\theta_e),\qquad
e=Tz,\qquad z_4=\eta(\theta-\theta_e).
\end{equation}
Define $\alpha=(G_\ell-LV\kappa)/I_\ell$. The reduced model is
\begin{equation}
\dot z=A_cz+B_c\delta F.
\end{equation}
\begin{equation}
A_c=T^TA_eT=
\begin{bmatrix}
-B_RR^2/I_\ell&-B_RR/I_\ell&-gm_r/I_\ell&\alpha/\eta&0\\
-B_RR/m_r&-B_R/m_r&0&(-g+RV\kappa)/\eta&0\\
R&1&0&0&0\\
\eta&0&0&0&0\\
0&0&0&-\kappa/\eta&0
\end{bmatrix},\qquad
B_c=\begin{bmatrix}R/I_\ell\\1/m_r\\0\\0\\0\end{bmatrix}.
\end{equation}
If compatibility is not imposed, write $e=Tz+w_e\zeta$, where
$\zeta=w_e^Te$ is constant. The projected equation then also contains
$T^TA_ew_e\zeta$; dropping this term does not remove the tracking mismatch.

\subsection{LQR and explicit feedback force}

Use the original lateral cost in the stated six-state order:
\begin{equation}
Q_{\rm lat}=\operatorname{diag}(2,2,20,20,1,0.5),\qquad \rho_F=0.1.
\end{equation}
For the analytic orthonormal basis above,
\begin{equation}
\boxed{Q_c=T^TQ_{\rm lat}T
=\operatorname{diag}\left(2,2,20,
\frac{20+\kappa^2}{1+\kappa^2},0.5\right).}
\end{equation}
Minimize $\int_0^\infty(z^TQ_cz+\rho_F\delta F^2)\,dt$.
Compute the continuous-time LQR gain for the reduced model:
\begin{equation}
K_c=\operatorname{lqr}(A_c,B_c,Q_c,\rho_F)
=(k_1,k_2,k_3,k_4,k_5).
\end{equation}
The lifted gain is explicitly
\begin{equation}
K_{\rm lat}=K_cT^T
=\begin{bmatrix}k_1&k_2&k_3&k_4/\eta&-\kappa k_4/\eta&k_5\end{bmatrix}.
\end{equation}
Consequently,
\begin{equation}
\boxed{\begin{aligned}
F(t)={}&m_r(g\theta_e+RV\Omega_e)
-k_1\sigma_1-k_2\sigma_r-k_3(r-r_e)\\
&-\frac{k_4}{\eta}
\bigl[(\theta-\theta_e)-\kappa(\sigma_3-\Omega_e)\bigr]
-k_5[\psi-\psi_{e0}-\Omega_e t].
\end{aligned}}
\end{equation}
On the compatible subspace this simplifies to
\begin{equation}
\boxed{\begin{aligned}
F(t)={}&m_r(g\theta_e+RV\Omega_e)
-k_1\sigma_1-k_2\sigma_r\\
&-k_3\left(r-\frac{G_T\theta_e+HV\Omega_e}{gm_r}\right)
-\eta k_4(\theta-\theta_e)
-k_5[\psi-\psi_{e0}-\Omega_e t].
\end{aligned}}
\end{equation}
The gains $k_j$ are computed numerically from LQR.


\section{Balancing control without $\psi$ feedback}
\label{sec:balance-no-heading}

For balancing only, omit the heading $\psi$ from the feedback state.

\subsection{Five-state model and controllability}
Define
\begin{equation}
y_b=(\sigma_1,\sigma_r,r,\theta,\sigma_3)^T,\qquad
\dot y_b=A_b y_b+B_bF,
\end{equation}
where
\begin{equation}
A_b=\begin{bmatrix}
-B_RR^2/I_\ell&-B_RR/I_\ell&-gm_r/I_\ell&G_\ell/I_\ell&LV/I_\ell\\
-B_RR/m_r&-B_R/m_r&0&-g&-RV\\
R&1&0&0&0\\
1&0&0&0&0\\
-\kappa&0&0&0&0
\end{bmatrix},\qquad
B_b=\begin{bmatrix}R/I_\ell\\1/m_r\\0\\0\\0\end{bmatrix}.
\end{equation}

\begin{equation}
w_b=\frac{1}{\eta}(0,0,0,\kappa,1)^T,\qquad
w_b^TA_b=0,\qquad w_b^TB_b=0.
\end{equation}
The controllability matrix still has rank four.

\subsection{Compatible balancing target}
For a prescribed lean $\theta_e$, write the initial and candidate target
states explicitly:
\begin{equation}
y_{b0}=\begin{bmatrix}
\sigma_1(0)\\\sigma_r(0)\\r(0)\\\theta(0)\\\sigma_3(0)
\end{bmatrix},\qquad
y_{be}(\Omega)=\begin{bmatrix}
0\\0\\(G_T\theta_e+HV\Omega)/(gm_r)\\\theta_e\\\Omega
\end{bmatrix}.
\end{equation}
Use the same left-nullspace target-selection condition:
\begin{equation}
\begin{aligned}
0=h_b(\Omega)
&=w_b^T[y_{b0}-y_{be}(\Omega)]\\
&=\frac{\kappa[\theta(0)-\theta_e]+\sigma_3(0)-\Omega}{\eta}.
\end{aligned}
\end{equation}
Thus
\begin{equation}
\begin{aligned}
\Omega_e&=\sigma_3(0)+\kappa[\theta(0)-\theta_e],\\
r_e&=\frac{G_T\theta_e+HV\Omega_e}{gm_r},\\
F_e&=m_r(g\theta_e+RV\Omega_e).
\end{aligned}
\end{equation}


\subsection{Analytic four-state reduction}
Let $e_b=y_b-y_{be}(\Omega_e)$ and $\delta F=F-F_e$.
An analytic orthonormal basis of $w_b^\perp$ is
\begin{equation}
T_b=\begin{bmatrix}
1&0&0&0\\0&1&0&0\\0&0&1&0\\
0&0&0&1/\eta\\0&0&0&-\kappa/\eta
\end{bmatrix},\qquad
T_b^TT_b=I_4,\qquad T_bT_b^T=I_5-w_bw_b^T.
\end{equation}

Transformation though $T_b$ gives the four-state reduced model with full rank controllability.

\subsection{Balancing LQR and force input}

After transformation, the reduced model is
\begin{equation}
\dot z_b=A_{bc}z_b+B_{bc}\delta F,\qquad
A_{bc}=T_b^TA_bT_b,\qquad B_{bc}=T_b^TB_b.
\end{equation}
The reduced cost is
\begin{equation}
Q_{bc}=T_b^TQ_bT_b=\operatorname{diag}\left(2,2,20,\frac{20+\kappa^2}{1+\kappa^2}\right),\qquad
Q_b=\operatorname{diag}(2,2,20,20,1).
\end{equation}

Recompute the four-state LQR gain:
\begin{equation}
K_{bc}=\operatorname{lqr}(A_{bc},B_{bc},Q_{bc},\rho_F)
=(\ell_1,\ell_2,\ell_3,\ell_4).
\end{equation}

\begin{equation}
\begin{aligned}
F(t)=F_e-K_{bc}z_b
={}&F_e-\ell_1\sigma_1-\ell_2\sigma_r-\ell_3(r-r_e)\\
&-\frac{\ell_4}{\eta}
[(\theta-\theta_e)-\kappa(\sigma_3-\Omega_e)].
\end{aligned}
\end{equation}
Replace the terms, it becomes
\begin{equation}
\begin{aligned}
F(t)={}&m_r(g\theta_e+RV\Omega_e)
-\ell_1\sigma_1-\ell_2\sigma_r\\
&-\ell_3\left(r-\frac{G_T\theta_e+HV\Omega_e}{gm_r}\right)
-\eta\ell_4(\theta-\theta_e).
\end{aligned}
\end{equation}



\end{document}


